\documentclass[12pt]{article}

% Packages
\usepackage[margin=1in]{geometry}
\usepackage{setspace}
\usepackage{titlesec}
\usepackage{hyperref}
\usepackage{graphicx}
\usepackage{enumitem}
\usepackage{parskip}
\usepackage{booktabs}
\usepackage{ulem}
\usepackage{fancyhdr}
\usepackage{tocloft}
\usepackage{float}
\usepackage{needspace}

\hypersetup{
    colorlinks=true,
    linkcolor=blue,
    urlcolor=blue
}

% Header/Footer
\pagestyle{fancy}
\fancyhf{}
\rhead{Cycle 3 Written Report}
\lhead{COMP 4710 Senior Design}
\cfoot{\thepage}

\begin{document}

%----------------------------
% Title Page
%----------------------------
\begin{titlepage}
    \centering
    \vspace*{2cm}
    {\LARGE \textbf{Cycle 3 Written Report}\\[1em]}
    {\large \textbf{Authors:}\\[0.5em]}
    {\large
        Trey Wendell\\
        Nathan Ferguson\\
        Patrick Bozeki\\
        Elizabeth Floyd\\
        Alan Leonard\\[1.5em]
    }
    {\large COMP 4710 Senior Design\\[0.5em]}
    {\large Department of Computer Science and Software Engineering\\
    Samuel Ginn College of Engineering, Auburn University\\}
    \vfill
\end{titlepage}

%----------------------------
% Table of Contents
%----------------------------
\tableofcontents
\newpage

%----------------------------
% 1. Executive Summary
%----------------------------
\section{System Metaphor (Executive Summary)}

Our system allows users to upload photos and transform them into buildable
LEGO\textsuperscript{\textregistered}-style brick mosaics and 3D sculptures,
complete with step-by-step building instructions and a full parts list.
Cycle 3 was the most feature-intensive cycle to date, delivering the complete
mosaic generation pipeline, an animated build guide, an AI-powered 3D mosaic
viewer, and a brand new 3D Model Builder that reconstructs a full LEGO sculpture
from a user-uploaded orbit video or multi-angle photo set.

The web application is built with Next.js (App Router), React, Tailwind CSS,
Three.js for 3D rendering, and a Python FastAPI service that hosts the machine
learning models. The entire stack runs in Docker Compose for local development
and is deployed to AWS Amplify (frontend) with a separate container host for the
ML service.

%----------------------------
% 2. Sprint Goals
%----------------------------
\section{Sprint Goals}

The goals established at the start of Cycle 3 were:

\begin{enumerate}[leftmargin=*]
    \item Implement end-to-end mosaic generation with color dithering and a live
          animated preview.
    \item Build a detailed result page with an interactive parts list, a
          before/after photo comparison, and step-by-step build guide
          instructions.
    \item Integrate a real-time 3D WebGL viewer that uses AI depth estimation to
          sculpt the mosaic into a bas-relief LEGO plaque.
    \item Deliver a 3D Model Builder that accepts user video or photos and
          produces a downloadable GLB model of a full LEGO sculpture.
    \item Containerize the ML service and resolve all blocking CI/CD issues to
          enable a live AWS Amplify deployment.
\end{enumerate}

All five goals were completed by the end of the cycle.

%----------------------------
% 3. System Architecture
%----------------------------
\section{System Architecture}

The application is organized as a monorepo with two primary services:

\begin{center}
\begin{tabular}{lll}
\toprule
\textbf{Service} & \textbf{Technology} & \textbf{Port} \\
\midrule
Frontend  & Next.js 16, React 19, Tailwind v4, Three.js & 3000 \\
ML API    & Python 3.11, FastAPI, Depth Anything V2      & 8000 \\
Database  & PostgreSQL (reserved)                         & 5432 \\
Cache     & Redis (reserved)                              & 6379 \\
\bottomrule
\end{tabular}
\end{center}

Docker Compose orchestrates all four services. The ML service ships in two
variants: a lightweight \texttt{Dockerfile.fastapi} (depth endpoint only,
suitable for free-tier hosts) and a \texttt{Dockerfile.pipeline} that includes
COLMAP, Open3D, and all geometry dependencies needed for full 3D reconstruction.

All generative routes are protected by Clerk authentication
(\texttt{/create}, \texttt{/build/*}, \texttt{/my-builds},
\texttt{/create-3d}, \texttt{/model/*}).

%----------------------------
% 4. Features Implemented
%----------------------------
\section{Features Implemented}

\subsection{Mosaic Generation (\texttt{/create} page)}

Users can upload up to five photos on the create page. The system quantizes
each image to a fixed 22-color LEGO palette using Floyd-Steinberg dithering,
which spreads color error to neighboring pixels and produces smooth tonal
gradients that closely approximate the original photograph. Three detail levels
control the resolution of the output mosaic:

\begin{center}
\begin{tabular}{lc}
\toprule
\textbf{Detail Level} & \textbf{Long-edge (studs)} \\
\midrule
Low    & 48 \\
Medium & 72 \\
High   & 104 \\
\bottomrule
\end{tabular}
\end{center}

A live animated preview plays immediately after upload: a glowing orange
scan-line sweeps across the canvas as the mosaic is drawn in, and re-plays
whenever the user switches detail levels. The finished build is saved to
browser \texttt{localStorage} along with a custom name and routed to the
result page.

\subsection{Build Result Page (\texttt{/build/[id]})}

The result page combines several interactive features into a single view:

\begin{itemize}[leftmargin=*]
    \item \textbf{Interactive parts list with hover highlighting.} Hovering any
          color row in the parts list dims all other bricks in the canvas,
          making it straightforward to locate a specific color group while
          building.

    \item \textbf{Before/After toggle.} A single button swaps between the
          original uploaded photo and the completed mosaic so users can judge
          the color translation at a glance.

    \item \textbf{Build Guide mode.} A step-through instruction mode modelled
          after official LEGO instructions. Each step isolates exactly one
          color; a progress bar tracks how many steps remain; and users can
          jump to any step by clicking the corresponding row in the parts list.

    \item \textbf{Export.} Users can download the mosaic as a PNG image or as
          a \texttt{.glb} 3D model file compatible with Blender and any
          WebGL-capable viewer.
\end{itemize}

\subsection{3D Mosaic Viewer with AI Depth Sculpting}

Clicking the 3D View button on the result page opens a full Three.js WebGL
scene. Each brick is rendered as a beveled box geometry with a cylindrical
stud on top; \texttt{THREE.InstancedMesh} is used throughout so the scene
remains interactive even for large mosaics with tens of thousands of bricks.
Orbit controls with damping let the user inspect the model from any angle, and
auto-rotate plays on load. Hover highlighting from the parts list is
synchronized into the 3D scene in real time by updating instance colors on the
\texttt{InstancedMesh}.

The central AI feature is \textbf{bas-relief depth sculpting}. When 3D View
opens, the frontend encodes the original photo as base64 and sends it to the
ML service:

\begin{verbatim}
POST /api/depth-grid
{ "image_b64": "...", "grid_w": 72, "grid_h": 54, "max_height": 8 }
\end{verbatim}

The service runs \textbf{Depth Anything V2 Small} (a state-of-the-art
monocular depth estimation model) to produce a per-cell height map, which is
returned to the frontend. Each brick column is then stacked to the
corresponding height, sculpting the flat mosaic into a 3D bas-relief plaque.
If the ML service is unavailable the viewer falls back to a single-layer flat
mosaic with a status indicator, so the page never breaks.

\subsection{3D Model Builder (\texttt{/create-3d} page)}

The 3D Model Builder allows users to upload an orbit video of any physical
object (10--30 seconds) or a ZIP of multi-angle photos (30--60 images). The
pipeline then reconstructs a full 3D LEGO sculpture and delivers it as a
downloadable GLB file.

\textbf{Frontend workflow:} After the user drops a file into the drag-and-drop
zone and clicks \textit{Start 3D Reconstruction}, the file is uploaded to the
ML service via \texttt{POST /api/ml/jobs}. A job ID is returned immediately and
the page begins polling:

\begin{verbatim}
GET /api/ml/jobs/${jobId}    // every 3 seconds
\end{verbatim}

A real-time elapsed timer increments every second so the user can gauge
progress. Full reconstruction typically takes 5--15 minutes depending on input
length. When the job status changes to \texttt{done}, the user is automatically
redirected to \texttt{/model/[jobId]}, where the finished GLB is rendered in
the \texttt{GlbViewer} component (orbit controls, auto-rotate, ACES filmic
tone-mapping).

\textbf{ML service pipeline:} Jobs are managed in an in-process thread pool
(\texttt{ThreadPoolExecutor}, 2 workers). Each job transitions through states
\texttt{PENDING} $\to$ \texttt{RUNNING} $\to$ \texttt{DONE} (or
\texttt{FAILED}). The full reconstruction pipeline---sparse reconstruction with
COLMAP, dense fusion with Open3D TSDF, brickification, and GLB export---runs in
a background thread imported from the \texttt{ptb\_ml} package. The default
lightweight container returns a descriptive error rather than crashing if the
pipeline dependencies are not present.

\textbf{API routes} proxying to the ML service:

\begin{center}
\begin{tabular}{ll}
\toprule
\textbf{Route} & \textbf{Purpose} \\
\midrule
\texttt{POST /api/ml/jobs}             & Upload file, start job, return \texttt{job\_id} \\
\texttt{GET  /api/ml/jobs/[jobId]}     & Poll status and progress string \\
\texttt{GET  /api/ml/jobs/[jobId]/glb} & Stream finished GLB to browser \\
\bottomrule
\end{tabular}
\end{center}

\subsection{Standalone Demo Script}

A command-line script (\texttt{ml/scripts/generate\_demo.py}) was added so
developers and evaluators can generate a GLB from a single image without
running Docker or the web application. It uses the same Floyd-Steinberg
quantization and Depth Anything V2 depth pipeline as the web app:

\begin{verbatim}
python3 scripts/generate_demo.py dog.jpg demo.glb --size 96 --max-height 12
\end{verbatim}

A \texttt{--flat} flag skips depth estimation for a quick flat-mosaic plaque.
The script prints a parts-list summary (color name and brick count per color)
to the terminal on completion. A sample \texttt{dog.jpg} and a pre-generated
\texttt{demo.glb} were committed alongside the script.

%----------------------------
% 5. Testing
%----------------------------
\section{Testing}

\begin{itemize}[leftmargin=*]
    \item \textbf{Mosaic correctness.} The dithering output was verified visually
          across multiple reference photos at all three detail levels, confirming
          that color transitions are smooth and the output closely matches the
          original image's palette.

    \item \textbf{3D viewer performance.} Mosaics at the High detail level
          (104-stud long edge) with depth sculpting enabled were profiled in
          Chrome DevTools; the scene maintained interactive frame rates using
          \texttt{InstancedMesh}.

    \item \textbf{ML endpoint.} The \texttt{/api/depth-grid} endpoint was tested
          end-to-end via the web app on several photos, and the generated
          height maps were verified to produce plausible bas-relief geometry in
          the 3D viewer.

    \item \textbf{3D pipeline.} The standalone \texttt{generate\_demo.py} script
          was run against a sample image to confirm that the depth model loads,
          quantization runs, and a valid GLB is written to disk. The committed
          \texttt{demo.glb} serves as a reference artifact.

    \item \textbf{Docker builds.} \texttt{docker compose up --build} was tested
          from a clean checkout to confirm that all four services start and
          that the frontend is accessible at \texttt{localhost:3000}.
\end{itemize}

%----------------------------
% 6. Team Contributions
%----------------------------
\section{Team Contributions}

\begin{center}
\begin{tabular}{lp{9cm}}
\toprule
\textbf{Member} & \textbf{Contributions} \\
\midrule
Trey Wendell &
    3D Model Builder (\texttt{/create-3d} page, job API routes, GLB viewer
    component), FastAPI job queue, \texttt{Dockerfile.pipeline},
    standalone demo script, README, AWS Amplify deployment config \\[0.5em]
Nathan Ferguson &
    Depth Anything V2 integration, \texttt{/api/depth-grid} endpoint,
    depth-driven bas-relief sculpting in \texttt{MosaicViewer3D} \\[0.5em]
Patrick Bozeki &
    Three.js 3D mosaic viewer (\texttt{MosaicViewer3D} component),
    instanced brick geometry, orbit controls, hover highlight sync \\[0.5em]
Alan Leonard &
    Mosaic generation (\texttt{/create} page), scan-line animation,
    Floyd-Steinberg dithering, \texttt{/my-builds} page, Clerk
    middleware updates \\
\bottomrule
\end{tabular}
\end{center}

%----------------------------
% 7. Obstacles and Resolutions
%----------------------------
\section{Obstacles and Resolutions}

\needspace{4\baselineskip}
\subsection{Broken DSINE Submodule Blocking CI}

A stale Git submodule pointer at \texttt{ml/vendor/DSINE} caused every AWS
Amplify build to fail immediately. The pointer was removed from the repository,
unblocking CI without affecting runtime functionality because DSINE was not yet
consumed by any live code path.

\needspace{4\baselineskip}
\subsection{\texttt{.dockerignore} Excluding the Build Route}

The Next.js \texttt{build/} output directory pattern in
\texttt{.dockerignore} was inadvertently also excluding the
\texttt{app/build/} route directory, causing the result page to return 404
inside the Docker container. The ignore rule was made more specific so that
only the Next.js \texttt{.next/} build cache is excluded.

\needspace{4\baselineskip}
\subsection{Keeping the Default Image Lightweight}

The full COLMAP reconstruction pipeline adds several gigabytes of dependencies.
To keep the default Docker image deployable on free-tier container hosts, the
pipeline dependencies are only present in \texttt{Dockerfile.pipeline}. The
\texttt{\_run\_pipeline} function in \texttt{app.py} catches the resulting
\texttt{ImportError} gracefully and marks the job as \texttt{FAILED} with a
human-readable error message rather than crashing the server.

%----------------------------
% 8. Retrospective
%----------------------------
\section{Retrospective}

\subsection{What Went Well}

\begin{itemize}[leftmargin=*]
    \item The Three.js \texttt{InstancedMesh} approach scaled to large mosaics
          without performance issues, validating the architecture choice early.
    \item Separating the lightweight and full-pipeline Docker images kept
          deployment options flexible and avoided forcing all team members to
          download the large pipeline image for everyday development.
    \item The scan-line animation created a strong sense of visual feedback
          that received positive reactions during internal demos.
\end{itemize}

\subsection{What Could Be Improved}

\begin{itemize}[leftmargin=*]
    \item Build and mosaic data are currently persisted only in
          \texttt{localStorage}, which prevents cross-device access and makes
          the data fragile. Moving to the PostgreSQL service is the highest
          priority for Cycle 4.
    \item The ML job queue has no rate limiting or concurrency guard beyond
          the two-worker thread pool, which could become a problem under
          simultaneous users.
    \item Mobile layout was not addressed this cycle; all pages are designed
          for desktop viewports.
\end{itemize}

%----------------------------
% 9. Plans for Next Cycle
%----------------------------
\section{Plans for Next Cycle}

\begin{enumerate}[leftmargin=*]
    \item \textbf{Database persistence.} Migrate build storage from
          \texttt{localStorage} to PostgreSQL so builds are accessible across
          devices and users.
    \item \textbf{Redis job queue.} Use Redis to queue and rate-limit ML jobs,
          preventing server overload under concurrent requests.
    \item \textbf{Brick cost estimator.} Add a real-time cost estimate to
          the parts list using current LEGO market prices from the BrickLink
          API.
    \item \textbf{Improved depth model.} Evaluate upgrading from Depth
          Anything V2 Small to the Medium or Large variant for higher-quality
          bas-relief sculpting on complex scenes.
    \item \textbf{Mobile-responsive layout.} Audit and update all pages for
          usability on mobile screen sizes.
\end{enumerate}

%----------------------------
% 10. Meeting Minutes
%----------------------------
\newpage
\section{Meeting Minutes}

\subsection{Team Meeting --- April 7, 2026}

\textbf{Attendees:} Trey Wendell, Nathan Ferguson, Patrick Bozeki, Alan Leonard\\
\textbf{Location:} Google Meet (remote) \quad \textbf{Duration:} 60 min\\[0.5em]

\textbf{Discussion:}
\begin{itemize}[leftmargin=*]
    \item Reviewed Cycle 2 feedback and agreed on Cycle 3 feature priorities.
    \item Assigned mosaic generation, scan-line animation, and \texttt{/my-builds} page to Alan Leonard.
    \item Assigned Three.js 3D mosaic viewer with instanced brick geometry to Patrick Bozeki.
    \item Assigned Depth Anything V2 integration and \texttt{/api/depth-grid} endpoint to Nathan Ferguson.
    \item Assigned 3D Model Builder (\texttt{/create-3d}), FastAPI job queue, and AWS Amplify deployment to Trey Wendell.
    \item Agreed to use Docker Compose for local development with all four services.
\end{itemize}

\textbf{Action Items:}
\begin{itemize}[leftmargin=*]
    \item All members: verify local Docker Compose stack builds and runs cleanly.
    \item Trey: assess COLMAP pipeline feasibility and stand up \texttt{Dockerfile.pipeline}.
    \item Nathan: prototype the depth estimation endpoint and confirm the Depth Anything V2 model downloads correctly.
    \item Alan: scaffold \texttt{/create} page and wire up Floyd-Steinberg dithering library.
\end{itemize}

\vspace{1em}
\subsection{Team Meeting --- April 14, 2026}

\textbf{Attendees:} Trey Wendell, Nathan Ferguson, Patrick Bozeki, Alan Leonard\\
\textbf{Location:} Google Meet (remote) \quad \textbf{Duration:} 60 min\\[0.5em]

\textbf{Discussion:}
\begin{itemize}[leftmargin=*]
    \item Alan demoed working mosaic generation at all three detail levels (Low 48, Medium 72, High 104 studs); scan-line animation was functional.
    \item Nathan confirmed \texttt{POST /api/depth-grid} endpoint returning per-cell height maps from Depth Anything V2 Small.
    \item Patrick showed early Three.js scene with instanced brick bodies; stud geometry and orbit controls still in progress.
    \item Trey reported FastAPI job queue operational with \texttt{PENDING} $\to$ \texttt{RUNNING} $\to$ \texttt{DONE} state transitions; \texttt{/create-3d} polling loop wired up.
    \item Identified that \texttt{.dockerignore} was excluding the \texttt{app/build/} route; assigned fix to Trey.
\end{itemize}

\textbf{Action Items:}
\begin{itemize}[leftmargin=*]
    \item Patrick: add stud geometry, shadow maps, and hover highlight sync to 3D viewer.
    \item Nathan: integrate depth grid response into \texttt{MosaicViewer3D} to drive column heights.
    \item Alan: implement build guide step-through mode and before/after toggle on result page.
    \item Trey: fix \texttt{.dockerignore}, resolve broken DSINE submodule blocking Amplify CI, and add standalone demo script.
\end{itemize}

\vspace{1em}
\subsection{Team Meeting --- April 21, 2026}

\textbf{Attendees:} Trey Wendell, Nathan Ferguson, Patrick Bozeki, Alan Leonard\\
\textbf{Location:} Google Meet (remote) \quad \textbf{Duration:} 60 min\\[0.5em]

\textbf{Discussion:}
\begin{itemize}[leftmargin=*]
    \item Merged all feature branches into \texttt{main}; resolved integration conflicts between the 3D viewer and depth endpoint changes.
    \item Confirmed \texttt{.dockerignore} fix resolved the \texttt{/build/[id]} 404 errors inside the container.
    \item Confirmed removal of DSINE submodule pointer unblocked Amplify CI; build pipeline now passes.
    \item Patrick demoed completed \texttt{MosaicViewer3D} with stud geometry, shadow maps, auto-rotate, and real-time hover highlight sync from the parts list.
    \item Nathan confirmed bas-relief depth sculpting working end-to-end: depth grid from ML service drives per-column brick heights in the 3D viewer.
    \item End-to-end walkthrough of all Cycle 3 features confirmed all five sprint goals complete.
\end{itemize}

\textbf{Action Items:}
\begin{itemize}[leftmargin=*]
    \item Trey: finalize README with full architecture documentation and deployment guide.
    \item All members: review written report draft before the April 28 meeting.
\end{itemize}

\vspace{1em}
\subsection{Team Meeting --- April 28, 2026}

\textbf{Attendees:} Trey Wendell, Nathan Ferguson, Patrick Bozeki, Alan Leonard\\
\textbf{Location:} Google Meet (remote) \quad \textbf{Duration:} 60 min\\[0.5em]

\textbf{Discussion:}
\begin{itemize}[leftmargin=*]
    \item Reviewed Cycle 3 written report as a group; confirmed accuracy of all sections.
    \item Discussed Cycle 4 priorities: database persistence via PostgreSQL, Redis job queue, and brick cost estimator.
    \item Agreed that mobile-responsive layout should be addressed early in Cycle 4 rather than deferred again.
    \item Reviewed live Amplify deployment; confirmed frontend accessible and ML service reachable via \texttt{FASTAPI\_INTERNAL\_URL}.
\end{itemize}

\textbf{Action Items:}
\begin{itemize}[leftmargin=*]
    \item All members: finalize and submit Cycle 3 written report.
    \item Trey: open Cycle 4 planning tickets and set up PostgreSQL schema for build persistence.
    \item Nathan: benchmark Depth Anything V2 Medium vs.\ Small to assess quality improvement for bas-relief sculpting.
\end{itemize}

%----------------------------
% 11. Sponsor Approval
%----------------------------
\newpage
\section{Sponsor Approval}

The undersigned sponsor representative has reviewed the Cycle 3 deliverables
for the AIPurpixBrickLab / PictoBrick project and approves the work completed
during this cycle.

\vspace{2em}

\noindent
\begin{tabular}{p{7cm} p{5cm}}
\rule{6.5cm}{0.4pt} & \rule{4.5cm}{0.4pt} \\
Sponsor Representative Signature   & Date \\[2.5em]
\rule{6.5cm}{0.4pt} & \\
Printed Name \& Title              & \\
\end{tabular}

\vspace{3em}

\noindent
The undersigned faculty advisor has reviewed and approved this report.

\vspace{2em}

\noindent
\begin{tabular}{p{7cm} p{5cm}}
\rule{6.5cm}{0.4pt} & \rule{4.5cm}{0.4pt} \\
Faculty Advisor Signature          & Date \\[2.5em]
\rule{6.5cm}{0.4pt} & \\
Printed Name \& Title              & \\
\end{tabular}

\end{document}
